\begin{figure}[htbp!]
	\centering
	\resizebox{\linewidth}{!}{%
		\begin{tikzpicture}[
			x=1.45cm, y=1.45cm, font=\scriptsize,
			trave/.style={line width=1.8pt, line cap=round, line join=round, black},
			skel/.style={line width=1.1pt, black},
			forza/.style={-{Stealth[length=1.9mm]}, line width=0.8pt, black},
			reaz/.style={-{Stealth[length=1.9mm]}, line width=0.8pt, gray!60},
			carico/.style={-{Stealth[length=1.4mm]}, line width=0.5pt, black},
			bordo/.style={line width=0.5pt, black},
			terra/.style={line width=0.6pt, black},
			fibra/.style={dashed, gray!70, line width=0.5pt},
			diagfill/.style={fill=none, draw=black, line width=0.8pt},
			nodoV/.style={circle, fill=white, draw=black, inner sep=0pt, minimum size=4.5pt, line width=0.8pt},
			segno/.style={circle, draw=black, fill=white, inner sep=0.4pt, line width=0.5pt, font=\tiny},
			nodo/.style={circle, fill=black, inner sep=0pt, minimum size=2.6pt},
			plab/.style={font=\footnotesize\bfseries},
			]
			\def\sT{0.65}\def\sTb{0.90} \def\sTc{1.30}  \def\sN{0.44}\def\sM{0.72}  % \sTb solo per ED \sTc solo per CD
			
			\newcommand{\setnodi}{%
				\coordinate (A) at (0,0); \coordinate (B) at (1.5,0);
				\coordinate (C) at (0.5,1); \coordinate (D) at (1.5,1); \coordinate (E) at (0.5,2);
				\coordinate (Mce) at (0.5,1.5); \coordinate (Mcd) at (1.0,1);}
			\newcommand{\cerniera}[1]{%
				\draw[terra] (#1)--++(-0.12,-0.20)--++(0.24,0)--cycle;
				\draw[terra] ($(#1)+(-0.17,-0.20)$)--($(#1)+(0.17,-0.20)$);
				\foreach \tk in {-0.14,-0.07,0,0.07,0.14}\draw[terra] ($(#1)+(\tk,-0.20)$)--++(-0.05,-0.08);}
			% struttura di riferimento (assi + cerniere interne + vincoli) da disegnare SOPRA i diagrammi
			\newcommand{\strutturaref}{%
				\draw[trave] (A)--(C); \draw[trave] (C)--(E)--(D); \draw[trave] (C)--(D)--(B);
				\cerniera{A}\cerniera{B}
				\node[nodo] at (D){};
				\node[nodoV] at (A){};\node[nodoV] at (B){};\node[nodoV] at (C){};\node[nodoV] at (E){};
				\coordinate (Dh) at ($(D)!0.06!(E)$);\node[nodoV] at (Dh){};}
			
			\newcommand{\diagLin}[7]{%
				\pgfmathsetmacro{\ax}{#5*#7*#3}\pgfmathsetmacro{\ay}{#5*#7*#4}%
				\pgfmathsetmacro{\bx}{#6*#7*#3}\pgfmathsetmacro{\by}{#6*#7*#4}%
				\filldraw[diagfill] (#1)--($(#1)+(\ax,\ay)$)--($(#2)+(\bx,\by)$)--(#2)--cycle;}
			\newcommand{\diagPar}[8]{%
				\pgfmathsetmacro{\ca}{2*(#5-2*#6+#7)}\pgfmathsetmacro{\cb}{-3*#5+4*#6-#7}\pgfmathsetmacro{\cc}{#5}%
				\foreach \k in {0,1,...,20}{%
					\pgfmathsetmacro{\tt}{\k/20}%
					\pgfmathsetmacro{\mv}{(\ca*\tt*\tt+\cb*\tt+\cc)*#8}%
					\pgfmathsetmacro{\ox}{\mv*#3}\pgfmathsetmacro{\oy}{\mv*#4}%
					\coordinate (pp\k) at ($(#1)!\tt!(#2)+(\ox,\oy)$);}%
				\filldraw[diagfill] (#1) foreach \k in {0,1,...,20}{ -- (pp\k)} -- (#2) -- cycle;}
			
			\newcommand{\hatchLin}[8]{% P1 P2 nx ny v1 v2 scala Nticks ; trattini grigi lungo la normale
				\foreach \hi in {0,1,...,#8}{%
					\pgfmathsetmacro{\ht}{\hi/#8}%
					\pgfmathsetmacro{\hv}{(#5+(#6-#5)*\ht)*#7}%
					\pgfmathsetmacro{\hox}{\hv*#3}\pgfmathsetmacro{\hoy}{\hv*#4}%
					\draw[gray,line width=0.4pt] ($(#1)!\ht!(#2)$)--($(#1)!\ht!(#2)+(\hox,\hoy)$);}}
			
			\newcommand{\hatchPar}[9]{% P1 P2 nx ny m0 mh m1 scala Nticks ; trattini grigi lungo la normale (parabola)
				\pgfmathsetmacro{\hca}{2*(#5-2*#6+#7)}\pgfmathsetmacro{\hcb}{-3*#5+4*#6-#7}\pgfmathsetmacro{\hcc}{#5}%
				\foreach \hi in {0,1,...,#9}{%
					\pgfmathsetmacro{\ht}{\hi/#9}%
					\pgfmathsetmacro{\hv}{(\hca*\ht*\ht+\hcb*\ht+\hcc)*#8}%
					\pgfmathsetmacro{\hox}{\hv*#3}\pgfmathsetmacro{\hoy}{\hv*#4}%
					\draw[gray,line width=0.3pt] ($(#1)!\ht!(#2)$)--($(#1)!\ht!(#2)+(\hox,\hoy)$);}}
			
			%==================== (a) STRUTTURA + CARICHI + FIBRE INFERIORI ====================
			\begin{scope}[shift={(0,3.5)}]
				\setnodi
				\draw[trave] (A)--(C); \draw[trave] (C)--(E)--(D); \draw[trave] (C)--(D)--(B);
				% fibre inferiori (solo tratteggio, lato intradosso, anche AC e DB)
				\draw[fibra] ($(A)+(0.085,-0.043)$)--($(C)+(0.085,-0.043)$);   % AC
				\draw[fibra] (0.5,0.90)--(1.5,0.90);                           % CD (sotto)
				\draw[fibra] (1.40,0)--(1.40,1);                               % DB (interno = sx)
				\draw[fibra] (0.60,1)--(0.60,2);                               % CE (interno = dx)
				\draw[fibra] ($(E)+(-0.07,-0.07)$)--($(D)+(-0.07,-0.07)$);     % ED (sotto)
				% coppie orarie
				\draw[black,line width=0.9pt,-{Stealth[length=1.8mm]}] (Mce)++(215:0.15) arc[start angle=215,end angle=-35,radius=0.15];
				\node[left] at ($(Mce)+(-0.16,0.02)$){$p\ell^2$};
				\draw[black,line width=0.9pt,-{Stealth[length=1.8mm]}] (Mcd)++(215:0.15) arc[start angle=215,end angle=-35,radius=0.15];
				\node[above] at ($(Mcd)+(-0.18,0.02)$){$p\ell^2$};
				% distribuito su ED
				\foreach \x in {0.5,0.625,...,1.501}\draw[carico] (\x,{2.5-\x+0.32})--(\x,{2.5-\x});
				\draw[bordo] (0.5,2.32)--(1.5,1.32);
				\node[above] at (1.0,1.92){$\tfrac{p}{\sqrt2}$};
				% orizzontale p su DB
				\foreach \y in {0,0.16,...,1.001}\draw[carico] (1.78,\y)--(1.5,\y);
				\draw[bordo] (1.78,0)--(1.78,1);
				\node[right] at (1.80,0.5){$p$};
				% reazioni grigie
				\draw[reaz] (A)--($(A)+(-0.45,0)$); \node[gray!70,below] at ($(A)+(-0.48,0.0)$){$\tfrac13 p\ell$};
				\draw[reaz] (A)--($(A)+(0,-0.55)$); \node[gray!70,left] at ($(A)+(0.5,-0.5)$){$\tfrac23 p\ell$};
				\draw[reaz] (B)--($(B)+(0.55,0)$); \node[gray!70,above] at ($(B)+(0.54,0)$){$\tfrac43 p\ell$};
				\draw[reaz] ($(B)+(0,-0.62)$)--($(B)+(0,-0.06)$); \node[gray!70,right] at ($(B)+(0,-0.46)$){$\tfrac53 p\ell$};
				\cerniera{A}\cerniera{B}
				\node[nodoV] at (A){}; \node[nodoV] at (B){}; \node[nodoV] at (C){}; \node[nodoV] at (E){};
				\coordinate (Dh) at ($(D)!0.06!(E)$); \node[nodoV] at (Dh){};
				\node at ($(A)+(-0.16,0.12)$){$A$}; \node at ($(B)+(-0.20,0.10)$){$B$};
				\node at ($(C)+(-0.16,0.04)$){$C$}; \node at ($(D)+(0.18,0.04)$){$D$}; \node at ($(E)+(-0.16,0.06)$){$E$};
				% pannello (a), in fondo al suo scope:
				\node[font=\tiny] at (0.75,-0.95){(a) struttura e forze};
			\end{scope}
			
			%==================== (b) TAGLIO T ====================
			\begin{scope}[shift={(3.7,3.5)}, segno/.append style={inner sep=0.2pt}]
				\setnodi
				\diagLin{C}{E}{-1}{0}{1}{1}{\sT}
				\diagLin{C}{D}{0}{-1}{0.16667}{0.16667}{\sTc}
				\diagLin{E}{D}{0.7071}{0.7071}{0.35355}{-0.35355}{\sTb}
				\diagLin{D}{B}{-1}{0}{-0.33333}{-1.33333}{\sT}
				\strutturaref
				
				\hatchLin{C}{E}{-1}{0}{1}{1}{\sT}{6}
				\diagLin{C}{E}{-1}{0}{1}{1}{\sT}
				\hatchLin{C}{D}{0}{-1}{0.16667}{0.16667}{\sTc}{6}
				\diagLin{C}{D}{0}{-1}{0.16667}{0.16667}{\sTc}
				\hatchLin{E}{D}{0.7071}{0.7071}{0.35355}{-0.35355}{\sTb}{9}
				\diagLin{E}{D}{0.7071}{0.7071}{0.35355}{-0.35355}{\sTb}
				\hatchLin{D}{B}{-1}{0}{-0.33333}{-1.33333}{\sT}{6}
				\diagLin{D}{B}{-1}{0}{-0.33333}{-1.33333}{\sT}
				
				% - nelle aste a taglio costante (centro della campitura)
				\node[segno,fill=white] at (0.5-0.5*\sT,1.5){$-$};                    % CE
				\node[segno,fill=white] at (1.0,1-0.5*0.16667*\sTc){$-$};              % CD
				\node[segno,fill=white] at ($(D)!0.5!(B)+(0.41667*\sT,0)$){$-$};      % DB (mezzo sbalzo medio)
				
				% ED: + quasi su E, - quasi su D
				\node[segno, fill=white] at ($(E)!0.08!(D)+(0.15*\sTb*0.7071,0.15*\sTb*0.7071)$){$+$};    % vicino a E
				\node[segno,fill=white] at ($(E)!0.92!(D)+(-0.15*\sTb*0.7071,-0.15*\sTb*0.7071)$){$-$};  % vicino a D
				
				\node[left,font=\tiny] at (0.5-1*\sT,1.5){$p\ell$};
				\node[below,font=\tiny] at (0.78,1-0.16667*\sTc){$\tfrac16 p\ell$};
				\node[above ,font=\tiny] at ($(E)+(0.35355*\sT*0.7071,0.35355*\sT*0.7071)$){$\tfrac{\sqrt2}{4}p\ell$};
				\node[below,font=\tiny] at ($(B)+(1.33333*\sT,0)$){$\tfrac43 p\ell$};
				\node[above,font=\tiny] at ($(D)+(0.33333*\sT,-0.05)$){$\tfrac13 p\ell$};
				\node[font=\tiny] at (0.75,-0.95){(b) taglio $T$};
			\end{scope}
			
			%==================== (c) NORMALE N ====================
			\begin{scope}[shift={(0,0)}]
				\setnodi
				\diagLin{A}{C}{-0.8944}{0.4472}{0.7454}{0.7454}{\sN}
				\diagLin{C}{E}{-1}{0}{0.5}{0.5}{\sN}
				\diagLin{E}{D}{-0.7071}{-0.7071}{-1.06066}{-1.76777}{\sN}
				\diagLin{C}{D}{0}{-1}{1.33333}{1.33333}{\sN}
				\diagLin{D}{B}{-1}{0}{-1.66667}{-1.66667}{\sN}
				\strutturaref
				
				\hatchLin{A}{C}{-0.8944}{0.4472}{0.7454}{0.7454}{\sN}{8}
				\hatchLin{C}{E}{-1}{0}{0.5}{0.5}{\sN}{6}
				\hatchLin{E}{D}{-0.7071}{-0.7071}{-1.06066}{-1.76777}{\sN}{9}
				\hatchLin{C}{D}{0}{-1}{1.33333}{1.33333}{\sN}{7}
				\hatchLin{D}{B}{-1}{0}{-1.66667}{-1.66667}{\sN}{6}
				
				% + nelle aste portanti a trazione
				\node[segno] at ($(A)!0.5!(C)+(0.5*0.7454*\sN*-0.8944, 0.5*0.7454*\sN*0.4472)$){$+$};  % AC
				\node[segno] at (0.5-0.5*0.5*\sN,1.5){$+$};                                            % CE
				\node[segno] at (1.0,1-0.5*1.33333*\sN){$+$};                                          % CD
				% - nelle aste compresse
				\node[segno] at ($(E)!0.5!(D)+(0.5*1.41422*\sN*0.7071,0.5*1.41422*\sN*0.7071)$){$-$};  % ED
				\node[segno] at ($(D)!0.5!(B)+(0.5*1.66667*\sN,0)$){$-$};                              % DB
				
				\node[above left,font=\tiny] at ($(A)!0.55!(C)+(0.7454*\sN*-0.8944,0.7454*\sN*0.4472)$){$\tfrac{\sqrt5}{3}p\ell$};
				\node[left,font=\tiny] at (0.5-0.5*\sN,1.82){$\tfrac12 p\ell$};
				
				% sforzo normale agli estremi di ED: -3√2/4 in E, -5√2/4 in D
				\node[left,font=\tiny, xshift=6pt] at ($(E)+(1.06066*\sN*0.7071,1.06066*\sN*0.7071)+(1.10,-0.1)$){$\tfrac{3\sqrt2}{4}p\ell$};
				\draw[-{Stealth[length=1.4mm]}, line width=0.4pt]($(E)+(1.06066*\sN*0.7071,1.06066*\sN*0.7071)+(0.58,-0.10)$) -- ($(E)+(1.06066*\sN*0.7071,1.06066*\sN*0.7071)$);
				
				\node[right,font=\tiny, xshift=-6pt, yshift=-8pt] at ($(D)+(1.76777*\sN*0.7071,1.76777*\sN*0.7071)+(0.16,0.02)$){$\tfrac{5\sqrt2}{4}p\ell$};
				\draw[-{Stealth[length=1.4mm]}, line width=0.4pt]
				($(D)+(1.76777*\sN*0.7071,1.76777*\sN*0.7071)+(0.08,-0.17)$) -- ($(D)+(0.88388*\sN*0.7071,0.88388*\sN*0.7071)$);
				
				%				\node[right,font=\tiny, yshift=-8pt] at ($(D)+(1.76777*\sN*0.7071,1.76777*\sN*0.7071)+(0.16,0.02)$){$\tfrac{5\sqrt2}{4}p\ell$};
				%				\draw[-{Stealth[length=1.4mm]}, line width=0.4pt]($(D)+(1.76777*\sN*0.7071,1.76777*\sN*0.7071)+(0.10,0.01)$) -- ($(D)+(0.88388*\sN*0.7071,0.88388*\sN*0.7071)$);
				%\node[above right,font=\tiny, xshift=22pt, yshift=-24pt] at ($(E)!0.55!(D)+(1.2*\sN*0.7071,1.2*\sN*0.7071)$){$\tfrac{5\sqrt2}{4}p\ell$};
				\node[below,font=\tiny] at (0.82,1-1.33333*\sN){$\tfrac43 p\ell$};
				\node[right,font=\tiny] at ($(D)!0.5!(B)+(1.66667*\sN,0)$){$\tfrac53 p\ell$};
				\node[font=\tiny] at (0.75,-0.75){(c) sforzo normale $N$};
			\end{scope}
			
			%==================== (d) MOMENTO M ====================
			\begin{scope}[shift={(3.7,0)}]
				\setnodi
				% CE (invariato): normale interna (1,0)
				\diagLin{C}{Mce}{-1}{0}{0}{0.5}{\sM}
				\diagLin{Mce}{E}{-1}{0}{-0.5}{0}{\sM}
				% CD ribaltato: normale (0,1)
				\diagLin{C}{Mcd}{0}{1}{0}{0.08333}{\sM}
				\diagLin{Mcd}{D}{0}{1}{-0.91667}{-0.83333}{\sM}
				% ED (invariato): normale interna (-1,-1)/sqrt2
				\diagPar{E}{D}{-0.7071}{-0.7071}{0}{0.125}{0}{\sM}
				% DB ribaltato: normale (-1,0)
				\diagPar{D}{B}{-1}{0}{0.83333}{0.54167}{0}{\sM}
				\strutturaref
				
				\hatchLin{C}{Mce}{-1}{0}{0}{0.5}{\sM}{3}
				\hatchLin{Mce}{E}{-1}{0}{-0.5}{0}{\sM}{3}
				\hatchLin{C}{Mcd}{0}{1}{0}{0.08333}{\sM}{3}
				\hatchLin{Mcd}{D}{0}{1}{-0.91667}{-0.83333}{\sM}{4}
				\hatchPar{E}{D}{-0.7071}{-0.7071}{0}{0.125}{0}{\sM}{9}
				\hatchPar{D}{B}{-1}{0}{0.83333}{0.54167}{0}{\sM}{6}
				
				
				% 1/2 pℓ²  → colmo del momento su CE (mezzeria)
				\draw[-{Stealth[length=1.4mm]}, line width=0.4pt]($(Mce)+(-0.5*\sM,0)+(-0.15,0.22)$) -- ($(Mce)+(0.35*\sM,0)$);
				\draw[-{Stealth[length=1.4mm]}, line width=0.4pt]($(Mce)+(-0.5*\sM,0)+(-0.15,0.22)$) -- ($(Mce)+(-0.25*\sM,0)$);
				\node[above,font=\tiny] at ($(Mce)+(-0.9*\sM,0.1)$){$\tfrac12 p\ell^2$};
				% +1/12 pℓ² : tratto a sinistra della coppia su CD (triangolo positivo sopra l'asse)
				\node[right,font=\tiny, xshift=-2pt] at ($(Mcd)+(0.60,0.34)$){$\tfrac{1}{12}p\ell^2$};
				\draw[-{Stealth[length=1.4mm]}, line width=0.4pt]($(Mcd)+(0.60,0.34)$) -- ($(Mcd)+(0,0.08333*\sM)$);
				% 11/12 pℓ²  → massimo su CD (a destra della coppia)
				\draw[-{Stealth[length=1.4mm]}, line width=0.4pt]($(Mcd)+(0,-0.45833*\sM)+(-0.4,-0.40)$) -- ($(Mcd)+(0,-0.65*\sM)$);
				\node[left,font=\tiny, xshift=-12pt] at ($(Mcd)!0.5!(D)+(0,-1.3*\sM)$){$\tfrac{11}{12}p\ell^2$};
				% 5/6 pℓ²  → momento al nodo rigido D (ramo CD, sotto D)
				\draw[-{Stealth[length=1.4mm]}, line width=0.4pt]($(D)+(0,-0.83333*\sM)+(0.28,0.08)$) -- ($(D)+(0,-0.83333*\sM)$);
				\node[above right,font=\tiny, xshift=14pt] at ($(D)+(-0.30,-0.83333*\sM)$){$\tfrac56 p\ell^2$};
				% archetto di continuità del momento al nodo rigido D (tra CD e DB)
				\draw[black, line width=0.4pt]($(D)+(-0.83333*\sM,0)$) arc[start angle=180, end angle=270, radius={0.83333*\sM}];
				% ordinata 1/8 pℓ² al colmo di ED, in nero
				\draw[black, line width=0.6pt]($(E)!0.5!(D)$) -- ($(E)!0.5!(D)+(-0.125*\sM*0.7071,-0.125*\sM*0.7071)$);
				\node[above right,font=\tiny, xshift=10pt, yshift=4pt] at ($(E)!0.55!(D)+(-0.125*\sM*0.7071,-0.125*\sM*0.7071)$){$\tfrac18 p\ell^2$};
				\draw[-{Stealth[length=1.4mm]}, line width=0.4pt]($(E)!0.5!(D)+(-0.125*\sM*0.7071,-0.125*\sM*0.7071)+(0.34,0.20)$)	--($(E)!0.5!(D)+(-0.125*\sM*0.7071,-0.125*\sM*0.7071)$);
				\node[font=\tiny] at (0.75,-0.75){(d) momento flettente $M$};
			\end{scope}
		\end{tikzpicture}
	}%
	\caption{Sistemi compositi a tre cerniere: (a)~schema strutturale
		con i carichi, le reazioni vincolari e le fibre inferiori di
		riferimento; (b)~taglio~$T$; (c)~sforzo normale~$N$;
		(d)~momento flettente~$M$. Sull'asta inclinata~$ED$ le ordinate
		sono tracciate ortogonalmente all'asse; al nodo rigido~$D$ il
		diagramma del momento è continuo dal lato interno. I tre
		diagrammi adottano fattori di scala indipendenti.}
	\label{fig:composito3c_NTM}
\end{figure}
